\section{Đặt vấn đề và tính cấp thiết của đề tài}
\label{sec:background}

Trong kỷ nguyên số hóa y tế và chương trình Chuyển đổi số Y tế Quốc gia, việc nâng cao chất lượng chẩn đoán, điều trị và đảm bảo an toàn người bệnh là mục tiêu hàng đầu tại các cơ sở khám chữa bệnh. Trong đó, khâu sử dụng thuốc (kê đơn, cấp phát, theo dõi tương tác và giám định bảo hiểm) là mắt xích phức tạp nhất, trực tiếp ảnh hưởng đến tính mạng bệnh nhân và hiệu quả kinh tế của bệnh viện.

Theo các nghiên cứu dịch tễ học lâm sàng của Tổ chức Y tế Thế giới (WHO), sự cố y khoa liên quan đến thuốc (Medication Errors) chiếm tỷ trọng lớn trong tổng số các biến cố bất lợi tại bệnh viện, trong đó phần lớn bắt nguồn từ sai sót trong quá trình chỉ định đơn thuốc (kê đơn quá liều, kê đơn có tương tác thuốc chống chỉ định, hoặc không hiệu chỉnh liều theo chức năng gan, thận). Tại Việt Nam, áp lực quá tải bệnh viện, số lượng bệnh nhân nội trú đông cùng với tỷ lệ người bệnh cao tuổi mắc đồng thời nhiều bệnh mạn tính (tăng huyết áp, đái tháo đường, suy tim, bệnh thận mạn) khiến tình trạng phối hợp đa thuốc (Polypharmacy) ngày càng phổ biến.

Mặt khác, công tác quản lý tài chính y tế đặt ra yêu cầu khắt khe về việc tuân thủ các quy định thanh toán Bảo hiểm Y tế (BHYT) theo Thông tư của Bộ Y tế. Sai sót trong việc kê đơn vượt tuyến kỹ thuật hoặc không đúng chỉ định chuyên môn dẫn đến tình trạng cơ quan bảo hiểm xuất toán hàng trăm triệu đồng, gây thiệt hại nghiêm trọng cho ngân sách cơ sở khám chữa bệnh.

Do đó, việc nghiên cứu và xây dựng một hệ thống phần mềm hỗ trợ ra quyết định lâm sàng (Clinical Decision Support System - CDSS) chuyên sâu về dược lâm sàng, kết hợp công nghệ trí tuệ nhân tạo có khả năng đối soát y văn pháp quy chuẩn hóa và cảnh báo tương tác thuốc thời gian thực là một nhu cầu cấp thiết, mang ý nghĩa khoa học và thực tiễn sâu sắc.
